\section{Pinned Contract Benchmark Construction}
\label{sec:benchmark_construction}

We construct a 100-file benchmark from WordPress~4.0 to evaluate migration under a pinned Rector contract.
The contract is a fixed Rector version, ruleset, and configuration.
We chose WordPress because it is large, long-lived, and exhibits substantial PHP version drift.

The benchmark is diagnostic by design.
We do not try to match the natural WordPress file distribution.
Instead, we optimize for coverage of the contract-induced obligation space.
Here, the obligation space is the set of Rector upgrade rule types that trigger in the corpus.
This choice makes results easier to interpret.
It also prevents rare but important obligations from being obscured by project-specific frequency effects.
The benchmark therefore measures progress against the pinned policy, not typicality of WordPress workloads.
This design preserves the full contract-induced rule space, enabling controlled and reproducible comparison of automated migration systems under a fixed upgrade policy.

The design follows standard guidance for controlled empirical studies and reproducible benchmark construction.\cite{wohlin2012experimentation}

\paragraph{Overview}
We build the benchmark in three steps:
\begin{enumerate}[leftmargin=*]
\item Run the pinned contract on the full corpus to obtain baseline triggers.
\item Characterize the triggered rule space (Table~\ref{tab:coverage_dimensions}).
\item Select 100 files using size stratification while preserving all rule types (Appendix~\ref{app:benchmark_algorithms}).
\end{enumerate}

We also define a 55-file execution slice.
We use it only as a loadability sentinel.
It flags new include-time failures under a fixed harness.
We interpret it only as a regression signal.

\begin{table}[htbp]
\centering
\caption{Coverage analysis dimensions used to characterize the triggered rule space under the pinned contract.}
\label{tab:coverage_dimensions}
\small
\begin{tabular}{@{}l p{0.50\columnwidth}@{}}
\toprule
\textbf{Dimension} & \textbf{Description} \\
\midrule
Rule type frequency & Corpus frequency of each Rector upgrade rule type. \\
PHP version mapping & Mapping rules to PHP~5.x, 7.x, and 8.x milestone families. \\
Rule categories & Distribution across AST constructs such as \texttt{FuncCall}, \texttt{Array\_}, and \texttt{Class\_}. \\
Size and complexity & Relationship between LOC and rule diversity. \\
Tail coverage & Preservation of infrequent rule types during sampling. \\
\bottomrule
\end{tabular}
\end{table}

\subsection{Inducing Contract Findings}
\label{sec:benchmark_static_analysis}

We analyze the full WordPress~4.0 corpus (482 PHP files) with the reference Rector configuration.
We pin \texttt{LevelSetList::UP\_TO\_PHP\_83} and the tool versions.
This keeps the contract stable across runs.
Static analysis triggers provide a consistent signal at scale.\cite{bessey2010billion,calcagno2015moving}

We record findings as file-level rule incidence.
A rule counts once per file if it triggers.
We use incidence because it is stable across repeated runs and configurations.
In-file counts are not consistently exposed by Rector.

We also record Rector dry-run patches.
They summarize edits suggested by the pinned contract.
We use them as a proxy for residual patch volume.\cite{falleri2014fine,fluri2007change}

Across the corpus we observe 1{,}331 file rule incidences spanning 48 upgrade rules.
This defines the obligation space used for benchmark construction and scoring.

\begin{table}[htbp]
\centering
\caption{Scoring criteria for size-stratified file selection.}
\label{tab:selection_criteria}
\small
\begin{tabular}{@{}l p{0.50\columnwidth}@{}}
\toprule
\textbf{Criterion} & \textbf{Role in File Selection} \\
\midrule
Rule diversity & Prefer files triggering many rule types. \\
Version change density & Reward stronger evidence of version-related changes. \\
Size suitability & Favor 100 to 2000 LOC and penalize very large files. \\
PHP family breadth & Prefer files spanning multiple rule families. \\
Key patterns & Extra weight for migration-relevant constructs. \\
\bottomrule
\end{tabular}
\end{table}

\subsection{Coverage Preserving Selection}
\label{sec:benchmark_selection}

We summarize baseline triggers to measure rule frequency, milestone coverage, and size effects.
Rare rule types matter, but uniform sampling can drop them.
We therefore enforce full rule coverage during selection.

Selection has two steps.
First, we stratify files by size and rank them using Table~\ref{tab:selection_criteria}.
Second, we enforce full rule coverage and trim to exactly 100 files while preserving coverage.
The procedure is deterministic.
We fix it before running any LLM experiments.

\begin{table}[htbp]
\centering
\caption{Target size distribution used for stratified initialization. The realized benchmark may differ after coverage enforcement.}
\label{tab:size_distribution}
\small
\begin{tabular}{@{}l c c@{}}
\toprule
\textbf{Category} & \textbf{LOC Range} & \textbf{Share} \\
\midrule
Small & 1--200 & 35\% \\
Medium & 201--500 & 30\% \\
Large & 501--1000 & 25\% \\
Extra large & 1001+ & 10\% \\
\bottomrule
\end{tabular}
\end{table}

\begin{table}[htbp]
\centering
\caption{Average migration complexity by file size in the released benchmark.}
\label{tab:migration_complexity}
\small
\begin{tabular}{@{}l c c c@{}}
\toprule
\textbf{Category} & \textbf{Files} & \textbf{Trig./File} & \textbf{Ver./File} \\
\midrule
Small & 31 & 3.2 & 2.9 \\
Medium & 31 & 5.4 & 4.8 \\
Large & 26 & 5.7 & 5.0 \\
Extra large & 12 & 8.9 & 6.2 \\
\bottomrule
\end{tabular}
\end{table}

\subsection{Benchmark Integrity}
\label{sec:benchmark_validation}

All selected files run through the pinned contract successfully.
The benchmark preserves all 48 upgrade rules.
It induces 521 file rule incidences spanning PHP~5.2 to 8.3.
Larger files trigger more rule types and span more milestones (Table~\ref{tab:migration_complexity}).
This matches prior links between structural complexity and transformation diversity.\cite{chidamber1994metrics,mccabe1976complexity}

Because coverage is the hard constraint, the final size mix can differ slightly from Table~\ref{tab:size_distribution}.

We also define a 55-file loadability slice.
We report only new failures (original pass, migrated fail) under the same harness.
We do not treat absolute pass rates as behavioral correctness.